\documentclass[12pt, a4paper]{article}

\usepackage[margin=1.8cm,top=2cm,bottom=2cm]{geometry}
\usepackage{amsmath,amssymb,amsfonts}
\usepackage{graphicx}
\usepackage[colorlinks=true,linkcolor=blue,citecolor=blue,urlcolor=blue]{hyperref}
\usepackage{bm,xcolor}
\usepackage[numbers,sort&compress]{natbib}
\usepackage{float}

\newcommand{\Geff}{G_{\mathrm{eff}}}
\newcommand{\meff}{m_{\mathrm{eff}}}
\newcommand{\heff}{\hbar_{\mathrm{eff}}}
\newcommand{\rhob}{\bar{\rho}}
\newcommand{\lP}{\ell_{\mathrm{P}}}

\begin{document}

\title{Environment-Dependent Vacuum Impedance:\\
The $\mu_0$--$\epsilon_0$ Splitting in the
Logarithmic Superfluid Vacuum}

\author{B. B. Kulangiev\\\small{Haifa, Israel}}
\date{April 2026}
\maketitle

\begin{abstract}
\noindent
The structural isomorphism between Maxwell's vacuum
constants ($\mu_0 \leftrightarrow \rho$,
$\epsilon_0 \leftrightarrow \kappa$) and the
mechanical properties of a compressible fluid is
well established. In the logarithmic superfluid
vacuum framework, this isomorphism acquires physical
content: $\mu_0$ and $\epsilon_0$ are not universal
constants but environment-dependent properties of the
vacuum condensate. We derive that near a gravitating
mass with dimensionless potential $U = GM/(c^2 r)$,
the vacuum permeability increases as
$\delta\mu_0/\mu_0 = +U$ while the permittivity
remains constant to first order:
$\delta\epsilon_0/\epsilon_0 = 0$. Consequently, the
speed of light decreases ($\delta c/c = -U/2$) while
the vacuum impedance increases
($\delta Z_0/Z_0 = +U/2$). This
$\mu_0$--$\epsilon_0$ splitting is a new prediction
that distinguishes the superfluid framework from
general relativity, which constrains only the
combined effect on $c$ through the PPN parameter
$\gamma$. We identify the fine structure constant
$\alpha$ as the hydrodynamic coupling efficiency of
a quantized vortex filament to the condensate
background, and show that the $\mu_0$--$\epsilon_0$
splitting predicts a gravitational environment
dependence of $\alpha$ at order $U^2$---below current
measurement precision but accessible to future
atomic clock comparisons in different gravitational
potentials.
\end{abstract}


% ====================================================================
\section{Introduction}
\label{sec:intro}
% ====================================================================

The impedance of free space,
\begin{equation}
    Z_0 = \sqrt{\frac{\mu_0}{\epsilon_0}}
    = \mu_0 c \approx 376.73\;\Omega\,,
\label{eq:Z0}
\end{equation}
is conventionally treated as a fixed property of
empty space. In transmission line theory, impedance
arises from the distributed inductance (inertia)
and capacitance (elasticity) of the
medium~\cite{jackson_1999}. For the vacuum, this
suggests that $\mu_0$ and $\epsilon_0$ encode
mechanical properties of a physical substrate.

The formal isomorphism between vacuum
electrodynamics and fluid
acoustics~\cite{landau_fluid} is well known:
\begin{equation}
    \mu_0 \leftrightarrow \rho\,,\qquad
    \epsilon_0 \leftrightarrow \kappa\,,\qquad
    c = \frac{1}{\sqrt{\mu_0\epsilon_0}}
    = \frac{1}{\sqrt{\rho\kappa}}\,,
\label{eq:isomorphism}
\end{equation}
where $\rho$ is the fluid density and
$\kappa = 1/(\rho c_s^2)$ is the adiabatic
compressibility. The acoustic impedance
$Z = \rho\,c_s = \sqrt{\rho/\kappa}$ maps to $Z_0$.

In the logarithmic superfluid vacuum
framework~\cite{kulangiev_gravity,kulangiev_geff},
this isomorphism is not merely formal. The vacuum is
a physical condensate with density $\rho_0$ and
sound speed $c_s$ determined by the logarithmic
equation of state~\cite{zloshchastiev_2023}. In a
companion paper~\cite{kulangiev_variable_c}, we
showed that $c_s$ varies across gravitational
environments as $\delta c/c = -U/2$, where
$U = GM/(c^2 r)$.

In this paper, we derive the \emph{individual}
environment dependence of $\mu_0$ and $\epsilon_0$
through the isomorphism. This yields a new
prediction---the $\mu_0$--$\epsilon_0$
splitting---that is absent in GR and testable in
principle.


% ====================================================================
\section{The Isomorphism in the Logarithmic Condensate}
\label{sec:isomorphism}
% ====================================================================

\subsection{The mapping}

The vacuum condensate obeys the relativistic
logarithmic Klein-Gordon
equation~\cite{zloshchastiev_2023,zloshchastiev_2020}.
Its Madelung decomposition yields a fluid with
density $\rho$ and sound
speed~\cite{kulangiev_gravity}:
\begin{equation}
    c_s^2 = \frac{c^2}{\ln(\rhob/\rho_c) + 2}\,.
\label{eq:cs}
\end{equation}
The adiabatic compressibility is:
\begin{equation}
    \kappa = \frac{1}{\rho\,c_s^2}\,.
\label{eq:kappa}
\end{equation}
Under the isomorphism~(\ref{eq:isomorphism}):
\begin{equation}
    \mu_0(\mathbf{x}) \propto \rho(\mathbf{x})\,,
    \qquad
    \epsilon_0(\mathbf{x}) \propto
    \kappa(\mathbf{x})
    = \frac{1}{\rho(\mathbf{x})\,c_s^2(\mathbf{x})}\,.
\label{eq:mapping}
\end{equation}
At the cosmological background
($\rho = \rho_\infty$, $c_s = c$), both reduce to
their standard values. Near a gravitating mass,
the density and sound speed shift, and so do
$\mu_0$ and $\epsilon_0$.

\subsection{The Painlev\'{e}-Gullstrand
environment}

From the PG flow
analysis~\cite{kulangiev_gravity,kulangiev_geff},
the vacuum density near a mass $M$ is:
\begin{equation}
    \rho(r) = \rho_\infty(1 + \epsilon)\,,\qquad
    \epsilon = U = \frac{GM}{c^2 r}\,.
\label{eq:rho_PG}
\end{equation}
The sound speed
is~\cite{kulangiev_variable_c}:
\begin{equation}
    c_s(r) \approx c\left(1 - \frac{\epsilon}{2}
    \right)\,.
\label{eq:cs_PG}
\end{equation}
The compressibility is:
\begin{equation}
    \kappa(r) = \frac{1}{\rho(r)\,c_s^2(r)}
    = \frac{1}{\rho_\infty(1+\epsilon)\,
    c^2(1-\epsilon)}
    \approx \frac{1}{\rho_\infty c^2}\,,
\label{eq:kappa_PG}
\end{equation}
where we used $(1+\epsilon)(1-\epsilon) \approx 1$
to first order in $\epsilon$.


% ====================================================================
\section{The \texorpdfstring{$\mu_0$--$\epsilon_0$}{μ0--ε0} Splitting}
\label{sec:splitting}
% ====================================================================

Collecting the first-order results:
\begin{align}
    \frac{\delta\mu_0}{\mu_0}
    &= \frac{\delta\rho}{\rho_\infty}
    = +U\,,
\label{eq:dmu}\\
    \frac{\delta\epsilon_0}{\epsilon_0}
    &= \frac{\delta\kappa}{\kappa_\infty}
    = 0 + O(U^2)\,,
\label{eq:deps}\\
    \frac{\delta c}{c}
    &= -\frac{1}{2}\left(
    \frac{\delta\mu_0}{\mu_0}
    + \frac{\delta\epsilon_0}{\epsilon_0}\right)
    = -\frac{U}{2}\,,
\label{eq:dc}\\
    \frac{\delta Z_0}{Z_0}
    &= \frac{1}{2}\left(
    \frac{\delta\mu_0}{\mu_0}
    - \frac{\delta\epsilon_0}{\epsilon_0}\right)
    = +\frac{U}{2}\,.
\label{eq:dZ}
\end{align}
The key result: near a gravitating mass, the vacuum
becomes magnetically denser ($\mu_0$ increases)
while remaining electrically unchanged ($\epsilon_0$
constant). The speed of light decreases because
$c = 1/\sqrt{\mu_0\epsilon_0}$, and the impedance
increases because
$Z_0 = \sqrt{\mu_0/\epsilon_0}$.

\begin{table}[H]
\centering
\caption{Environment dependence of vacuum
electromagnetic properties at the solar limb
($U \approx 2.1 \times 10^{-6}$).}
\begin{tabular}{lcc}
\hline
Quantity & Fractional shift & Value \\
\hline
$\delta\mu_0/\mu_0$ & $+U$ & $+2.1 \times 10^{-6}$ \\
$\delta\epsilon_0/\epsilon_0$ & $0$ & $0$ \\
$\delta c/c$ & $-U/2$ & $-1.1 \times 10^{-6}$ \\
$\delta Z_0/Z_0$ & $+U/2$ & $+1.1 \times 10^{-6}$ \\
\hline
\end{tabular}
\label{tab:splitting}
\end{table}

\subsection{Comparison with general relativity}

In GR, the coordinate speed of light near a mass
varies as $\delta c/c \sim -2U$ (in isotropic
coordinates), determined by the PPN parameter
$\gamma = 1$. However, GR does not decompose this
into separate $\mu_0$ and $\epsilon_0$ contributions
because GR treats the vacuum as geometric, not
material. The splitting~(\ref{eq:dmu})--(\ref{eq:deps})
is a prediction unique to the superfluid framework.

The Cassini constraint ($|\gamma - 1| < 2.3 \times
10^{-5}$~\cite{bertotti_2003}) constrains the
\emph{combined} effect on $c$, not the individual
$\mu_0$ and $\epsilon_0$ variations. The splitting
is consistent with Cassini.

\subsection{Physical interpretation}

The asymmetry has a clear physical origin. Near a
mass, the condensate density increases (more fluid
per unit volume), which increases the inductive
response ($\mu_0$). However, the compressibility
$\kappa = 1/(\rho c_s^2)$ stays constant because
the increase in $\rho$ is exactly compensated by the
decrease in $c_s^2$: a denser fluid is also stiffer,
and the two effects cancel in the elastic response.

In transmission line language: the inductance per
unit length increases near a mass, but the
capacitance per unit length stays the same. The line
becomes slower (lower $c$) and higher impedance
(higher $Z_0$).


% ====================================================================
\section{The Fine Structure Constant}
\label{sec:alpha}
% ====================================================================

The fine structure constant is related to the vacuum
impedance by~\cite{jackson_1999}:
\begin{equation}
    \alpha = \frac{e^2}{4\pi\epsilon_0\hbar c}
    = \frac{Z_0}{2R_K}\,,
\label{eq:alpha_Z}
\end{equation}
where $R_K = h/e^2$ is the von Klitzing constant.

In the condensate picture, $\alpha$ is the coupling
efficiency between a quantized vortex (the charged
particle) and the background condensate. The von
Klitzing constant $R_K$ is the impedance of a single
quantized circulation quantum; $\alpha$ measures how
much of this ideal circulation couples to the bulk
medium.

Since $Z_0$ is environment-dependent
($\delta Z_0/Z_0 = +U/2$), $\alpha$ is in principle
environment-dependent as well. However, $e$ and
$\hbar$ are topological quantities (the vortex
winding number and the phase quantum) and do not
depend on the local condensate density. Therefore:
\begin{equation}
    \frac{\delta\alpha}{\alpha}
    = \frac{\delta Z_0}{Z_0}
    - \frac{\delta R_K}{R_K}
    = \frac{U}{2} - 0 = \frac{U}{2}\,.
\label{eq:dalpha}
\end{equation}
Wait---this requires care. From the definition
$\alpha = e^2 Z_0/(4\pi\hbar)$:
if $Z_0$ increases, $\alpha$ increases. But
$\alpha = e^2/(4\pi\epsilon_0\hbar c)$, and since
$\epsilon_0$ is constant and $c$ decreases:
\begin{equation}
    \frac{\delta\alpha}{\alpha}
    = -\frac{\delta c}{c}
    = +\frac{U}{2}\,.
\label{eq:dalpha2}
\end{equation}
Both routes agree. At the solar limb:
$\delta\alpha/\alpha \sim 10^{-6}$. At the surface
of a white dwarf ($U \sim 10^{-4}$):
$\delta\alpha/\alpha \sim 5 \times 10^{-5}$.

Current atomic clock comparisons between different
gravitational potentials (e.g., ground vs.\ ISS)
constrain $\delta\alpha/\alpha < 10^{-7}$ per unit
$\Delta U$~\cite{safronova_2018}. The predicted
$\delta\alpha/\alpha = U/2$ is at the edge of
current sensitivity and may be testable with
next-generation optical lattice clocks.


% ====================================================================
\section{Observational Tests}
\label{sec:tests}
% ====================================================================

\subsection{Test 1: Atomic spectroscopy in different
gravitational potentials}

The transition frequencies of atoms depend on
$\alpha$ as $\nu \propto \alpha^q$, where $q$ is a
sensitivity coefficient that varies between
transitions~\cite{safronova_2018}. If
$\delta\alpha/\alpha = U/2$, transitions with
different $q$ values will shift by different amounts
in different gravitational potentials.

Comparing two transitions with $q_1 \neq q_2$ in
the same atom at two gravitational potentials
$U_1 \neq U_2$:
\begin{equation}
    \frac{\delta(\nu_1/\nu_2)}{\nu_1/\nu_2}
    = (q_1 - q_2)\,\frac{\delta\alpha}{\alpha}
    = (q_1 - q_2)\,\frac{\Delta U}{2}\,.
\label{eq:clock_test}
\end{equation}
For Yb$^+$ ($q_1 = -5.95$ for the octupole
transition) and Cs ($q_2 = 0.83$):
$q_1 - q_2 \approx -6.8$. Between Earth's surface
($U_\oplus \approx 7 \times 10^{-10}$) and a solar
probe at 10 solar radii ($U \approx 2 \times
10^{-7}$), $\Delta U \approx 2 \times 10^{-7}$,
giving:
\begin{equation}
    \frac{\delta(\nu_1/\nu_2)}{\nu_1/\nu_2}
    \approx 7 \times 10^{-7}\,.
\end{equation}
This is within the sensitivity of current optical
clock technology ($\sim 10^{-18}$ fractional
frequency), but the challenge is deploying clocks
at sufficiently different potentials.

\subsection{Test 2: Quasar absorption spectra}

Measurements of $\alpha$ from quasar absorption
systems at different redshifts have reported
tentative spatial variations at the $10^{-5}$
level~\cite{webb_2011}. If these systems reside in
different gravitational environments (voids vs.\
clusters), the superfluid framework predicts a
correlation between $\delta\alpha/\alpha$ and the
large-scale gravitational potential---the same
environment dependence predicted for the speed of
light~\cite{kulangiev_variable_c}.


% ====================================================================
\newpage
\section{Discussion}
\label{sec:discussion}
% ====================================================================

\subsection{What is new}

The hydrodynamic isomorphism $\mu_0 \leftrightarrow
\rho$, $\epsilon_0 \leftrightarrow \kappa$ is well
known~\cite{landau_fluid,jackson_1999}. What the
logarithmic superfluid framework adds is:

(i) A specific equation of state (Eq.~\ref{eq:cs})
that determines how $\rho$ and $c_s$ vary with the
gravitational potential.

(ii) The $\mu_0$--$\epsilon_0$ splitting: $\mu_0$
increases by $+U$ while $\epsilon_0$ stays constant.
This is a derived result, not postulated.

(iii) A prediction for the environment dependence of
$\alpha$: $\delta\alpha/\alpha = U/2$.

(iv) Two observational tests that can confirm or
falsify these predictions.

\subsection{Relation to Scharnhorst effect}

The Scharnhorst effect~\cite{scharnhorst_1990}
predicts a fractional speed increase between Casimir
plates of order $\alpha^2/(m_e a)^4$, arising from
QED vacuum polarization. In the superfluid framework,
the Casimir geometry modifies the local condensate
density (a boundary condition effect), producing a
speed shift through the equation of state. The two
mechanisms are complementary: the Scharnhorst effect
is a perturbative QED correction; the condensate
effect is a non-perturbative thermodynamic response.
A detailed comparison requires computing the
condensate density profile between conducting plates,
which we defer to future work.

\subsection{Limitations}

The proportionality constants in the isomorphism
($\mu_0 \propto \rho$) are dimensional and require
the gauged Lagrangian of the companion
paper~\cite{kulangiev_maxwell} to fix. The present
paper derives only the \emph{fractional} variations,
which are independent of the proportionality
constants.

The prediction $\delta\alpha/\alpha = U/2$ assumes
that $e$ and $\hbar$ are environment-independent.
This is justified if charge and action are
topological quantities (winding number and phase
quantum), but a rigorous proof requires deriving
$e$ and $\hbar$ from condensate parameters---an
open problem.


% ====================================================================
\section{Conclusion}
\label{sec:conclusion}
% ====================================================================

We have derived, within the logarithmic superfluid
vacuum framework, that the vacuum electromagnetic
constants $\mu_0$ and $\epsilon_0$ respond
asymmetrically to the gravitational environment:

(1) Near a gravitating mass ($U = GM/c^2 r$):
$\delta\mu_0/\mu_0 = +U$ (vacuum becomes
magnetically denser), while
$\delta\epsilon_0/\epsilon_0 = 0$ (electrically
unchanged).

(2) The speed of light decreases:
$\delta c/c = -U/2$.

(3) The vacuum impedance increases:
$\delta Z_0/Z_0 = +U/2$.

(4) The fine structure constant increases:
$\delta\alpha/\alpha = +U/2$.

The $\mu_0$--$\epsilon_0$ splitting is absent in
general relativity and constitutes a falsifiable
prediction of the superfluid framework, testable
through atomic clock comparisons in different
gravitational potentials or through correlations
between quasar absorption measurements of $\alpha$
and the large-scale gravitational environment.

% ====================================================================
\newpage
\section*{Acknowledgments}
% ====================================================================

Claude (Anthropic, claude-opus-4-6) and Gemini (Google, Gemini 3.1 Pro) were used as editorial and computational tools during manuscript preparation. All scientific content, equations, and
interpretations are the author's own.

\begin{thebibliography}{20}

\bibitem{kulangiev_gravity}
B.~B.~Kulangiev,
Conditions for Emergent Gravitational Light Bending from a Logarithmic Superfluid Vacuum,
Preprint (2026).
\url{https://kulangiev.github.io/#paper1}

\bibitem{kulangiev_geff}
B.~B.~Kulangiev,
Emergent Newtonian Gravity and the Gravitational Feynman-Onsager Relation in the Logarithmic Superfluid Vacuum,
Preprint (2026).
\url{https://kulangiev.github.io/#combined}

\bibitem{kulangiev_maxwell}
B.~B.~Kulangiev,
Emergent Electromagnetism from the Gauged Logarithmic Superfluid Vacuum,
Preprint (2026).
\url{https://kulangiev.github.io/#paper8}

\bibitem{kulangiev_variable_c}
B.~B.~Kulangiev,
A Testable Prediction of the Logarithmic Superfluid Vacuum: Environment-Dependent Speed of Light,
Preprint (2026).
\url{https://kulangiev.github.io/#paper9}

\bibitem{zloshchastiev_2023}
K.~G.~Zloshchastiev,
Derivation of emergent spacetime metric, gravitational
potential and speed of light in superfluid vacuum
theory,
Universe \textbf{9}, 234 (2023).

\bibitem{zloshchastiev_2020}
K.~G.~Zloshchastiev,
An alternative to dark matter and dark energy:
Scale-dependent gravity in superfluid vacuum theory,
Universe \textbf{6}, 180 (2020).

\bibitem{jackson_1999}
J.~D.~Jackson,
\emph{Classical Electrodynamics}, 3rd ed.\
(Wiley, 1999).

\bibitem{landau_fluid}
L.~D.~Landau and E.~M.~Lifshitz,
\emph{Fluid Mechanics}, 2nd ed.\
(Butterworth-Heinemann, 1987).

\bibitem{bertotti_2003}
B.~Bertotti, L.~Iess, and P.~Tortora,
A test of general relativity using radio links with
the Cassini spacecraft,
Nature \textbf{425}, 374 (2003).

\bibitem{scharnhorst_1990}
K.~Scharnhorst,
On propagation of light in the vacuum between plates,
Phys.\ Lett.\ B \textbf{236}, 354 (1990).

\bibitem{safronova_2018}
M.~S.~Safronova \emph{et al.},
Search for new physics with atoms and molecules,
Rev.\ Mod.\ Phys.\ \textbf{90}, 025008 (2018).

\bibitem{webb_2011}
J.~K.~Webb \emph{et al.},
Indications of a spatial variation of the fine
structure constant,
Phys.\ Rev.\ Lett.\ \textbf{107}, 191101 (2011).

\end{thebibliography}

\end{document}
